\section{The additive basis carries the CRT stratification}
\label{sec:crt}

The multiplicative basis is where the clock lives. The additive basis is not empty, and what
it holds is the stratification itself.

\subsection{The key additive frequencies are the CRT-component duals}
\label{sec:crt-law}

Let $n = \prod_{i=1}^{\omega(n)} q_{i}$ be the factorisation into \emph{maximal} prime powers
$q_{i} = p_{i}^{\,\alpha_{i}}$, so that the Chinese remainder theorem gives
$\mathbb{Z}/n\mathbb{Z} \cong \prod_{i} \mathbb{Z}/q_{i}\mathbb{Z}$. Define the predicted
additive frequency set
\begin{equation}
  \label{eq:crtset}
  P(n) \;=\; \bigcup_{i=1}^{\omega(n)}
  \Bigl\{\, t \cdot \tfrac{n}{q_{i}} \;:\; t \ge 1,\; t \cdot \tfrac{n}{q_{i}} \le \lfloor n/2 \rfloor \,\Bigr\}
  \;\subseteq\; \bigl\{1, \dots, \lfloor n/2 \rfloor \bigr\}.
\end{equation}
The ceiling at $\lfloor n/2 \rfloor$ is not cosmetic: the additive spectrum folds conjugate
frequencies, so it has exactly $\lfloor n/2 \rfloor$ bins and $k$ and $n - k$ are the same
frequency. Dropping the ceiling doubles $|P(n)|$ and the test no longer means anything. At
$n = 165$, \eqref{eq:crtset} gives $P = \{15, 30, 33, 45, 55, 60, 66, 75\}$; at $n = 120$ it
gives $\{15, 24, 30, 40, 45, 48, 60\}$.

\begin{quote}
The key additive frequencies of a grokked $a \cdot b \bmod n$ transformer are exactly $P(n)$:
one family per maximal prime power, the additive duals of the Chinese-remainder
components.
\end{quote}

The mechanism is elementary and needs no network. The indicator of the multiples of
$d$ has additive Fourier support on the multiples of $n/d$, and by \eqref{eq:jclass} each
$\mathcal{J}$-class is a union of arithmetic progressions generated by the CRT components.
At $\omega(n) = 1$ we have $n/q_{1} = 1$, so \eqref{eq:crtset} is every one of the
$\lfloor n/2 \rfloor$ bins: at $n = 113$, all $56$ of them. The prediction is then vacuous
and a prime power can never be counted as support.

Pre-registered before the moduli were run, and tested by a threshold-free permutation
enrichment over $20{,}000$ permutations. It holds at $9$ of $9$ testable moduli at $3$ of $3$
seeds each, every one at $p < 0.01$:

\begin{center}
\begin{tabular}{rcc}
  \toprule
  $n$ & enrichment, seeds 0/1/2 & seeds $p<0.01$ \\
  \midrule
  54 & 2.7/1.5/6.1 & 3/3 \\
  63 & 4.7/1.6/1.6 & 3/3 \\
  75 & 3.7/7.6/3.3 & 3/3 \\
  98 & 7.6/6.8/8.3 & 3/3 \\
  99 & 5.9/4.2/7.3 & 3/3 \\
  100 & 3.7/6.9/5.5 & 3/3 \\
  105 & 11.8/14.4/11.0 & 3/3 \\
  143 & 10.6/6.3/3.4 & 3/3 \\
  147 & 8.6/5.1/4.1 & 3/3 \\
  \bottomrule
\end{tabular}
\end{center}

% STE descriptive
\begin{figure}[t]
  \centering
  \includegraphics[width=\textwidth]{../figures/F6_crt.pdf}
  \caption[The predicted set is fixed by the factorisation before any model is trained, and
    the energy is on it.]{\textbf{The predicted set is fixed by the factorisation before any
    model is trained, and the energy is on it.} \emph{Top:} additive spectral energy of
    $W_{E}$ against frequency $k$, in raw residue coordinates, where $P(n)$ of
    \eqref{eq:crtset} lives. Bars in $P(n)$ are red, and the rest are grey. \emph{Bottom:}
    the threshold-free test. The grey histogram is the enrichment of $20{,}000$ random
    frequency sets of the same cardinality, and the red rule is the observed enrichment. The
    axis is logarithmic because the null concentrates near $1\times$.\par
    The rightmost column is \citet{chen2026multiplication}'s published checkpoint scored by
    our pipeline, which is the calibration for the other three. At $n = 119$ the effect is
    $3.30\times$ rather than $14.08\times$ and still lies beyond every draw. That modulus was
    enriched all along and invisible to a $5\times$-median threshold, which is why the test is
    threshold-free. Prime powers are omitted because $n/q_{1} = 1$ makes $P(n)$ every bin.
    That is vacuous by construction, not an absent measurement. The figure regenerates from
    \texttt{src/viz/plots.py::crt\_law}.\par
    The predicted set, the energy and the null come from \texttt{test\_crt\_law.py}. The
    generator asserts the published $14.08\times$ at $n = 165$ before it draws, because a
    bin-convention error here is silent at every modulus but one.}
  \label{fig:crt}
\end{figure}
% STE end

$8$ of those $12$ moduli (the nine of the table together with $165$, $120$ and $119$) are
non-square-free. Two further CRT-testable moduli were added
later: $n = 123 = 3 \cdot 41$ at enrichment $7.59 / 3.28 / 7.43$, $p < 5 \times 10^{-5}$ in $3/3$. At the
six original moduli the law holds at five seeds. $165$ reads $17.4, 13.3, 11.2, 11.2, 11.8$
and $120$ reads $13.8, 11.4, 9.8, 7.6, 13.3$, both $5/5$. $119$ reads
$3.3, 11.1, 5.2, 7.9, 6.4$, also $5/5$: it was enriched all along but invisible to a threshold
rather than absent.

The predicted and observed sets are \emph{exactly} equal, with no misses and no extras, at
$n = 165$ on both our own model and \citet{chen2026multiplication}'s published checkpoint, and
at $n = 120$. The second is the sharper test, because $120 = 2^{3} \cdot 3 \cdot 5$ has
generators $120/8$, $120/3$ and $120/5$: \emph{maximal prime powers, not primes}. Replacing
$q_{i}$ by $p_{i}$ in \eqref{eq:crtset} predicts $\{24, 40, 48, 60\}$ and misses
$\{15, 30, 45\}$, so the distinction is testable and the data decides it.

We asked the same question of the stratification with no network involved, and the answer
bounds the law rather than confirming it. The $\mathcal{J}$-class indicators of
$\mathbb{Z}/n\mathbb{Z}$ are sparse in the additive basis.
$\mathrm{Gini}_{\text{add}}$ is $0.406$ at $n = 121$, $0.450$ at $165$, $0.400$ at $120$, and
$\mathbf{0.000}$ at $n = 113$, a field having no stratification to encode. Their energy is
also strongly enriched on $P(n)$: $4.87\times$ at $n = 165$, $3.32\times$ at $120$ and
$7.15\times$ at $119$, every one beyond all $20{,}000$ permutations. But they do not
\emph{select} $P(n)$. Under the same detector and the same bin convention as the model
measurements, the indicators' key set equals $P(n)$ at only $2$ of the $10$ composite moduli
we tested. At $n = 165$ it holds $3$ of the $8$ predicted frequencies, and at $n = 120$ it is
$\{20, 40, 60\}$. Here $20 = 120/6$ is the dual of a divisor that is \emph{not} a maximal
prime power, so it lies outside $P(n)$ and outside the prime variant alike.

This is structural. The transform of $\mathbf{1}_{J_{d}}$ at frequency $k$ is the Ramanujan
sum $c_{q}(k) = \mu(q/g)\,\varphi(q)/\varphi(q/g)$, with $q = n/d$ and $g = \gcd(k, q)$. It
has magnitude at most $1$ wherever $k$ is coprime to $n$ and reaches $\varphi(q)$ where $k$
shares all of $q$. So the stratification is not confined to any divisor family: at
$n = 165$ its energy is nonzero at all $82$ frequencies. But it concentrates on the
frequencies that share a factor with $n$, the multiples of the duals of \emph{every}
divisor. At $n = 165$ those are $42$ frequencies, carrying $30.5\times$ the mean energy of
the other $40$, against $|P(165)| = 8$.

That makes the contrast with the network the result, and it runs in the paper's favour. On the same run, with one instrument: the model's key set is all $8$ of $P(165)$
and all $7$ of $P(120)$, with no misses and no extras, while the indicators select a subset
of the much larger divisor-dual family. Choosing the maximal-prime-power subfamily out of
the divisor duals is therefore something the trained network does, not something the algebra
hands it. The enrichment figures in this paragraph were computed after the set-equality test
they replace had failed, and are exploratory. The contrast table is not: both of its
columns use the detector fixed before the measurement.

The vacuous case excludes $113$, $121$, $125$, $127$, $131$, $128$, $49$, $81$ and $169$:
every prime and every prime power in our set.

\subsection{The early enrichment ramp is not the grokking circuit}
\label{sec:crt-ramp}

The CRT enrichment rises monotonically from early in training, it is specific to the CRT
structure against three size-matched nulls, and it would be a convenient progress measure. It
is not one, and this subsection is the negative result that says so.

We trained three seeds at $n = 63$ that did not generalise (test-accuracy window medians
$0.1801$, $0.1655$, $0.3183$, with training accuracy $1.000$ from step $2{,}000$, a clean
memorisation arm) and scored the same ramp, under criteria committed before the runs
existed:

\begin{center}
\begin{tabular}{lcccc}
  \toprule
  seed & final accuracy & enrichment at $1$k & enrichment at $20$k & ramp $\rho$ \\
  \midrule
  0 & $0.1803$ & $1.397$ & $\mathbf{1.653}$ & $\mathbf{1.000}$ \\
  1 & $0.1655$ & $1.339$ & $\mathbf{1.459}$ & $\mathbf{0.988}$ \\
  2 & $0.3177$ & $1.422$ & $\mathbf{1.655}$ & $\mathbf{0.995}$ \\
  \bottomrule
\end{tabular}
\end{center}

The failed runs show the same monotone ramp as grokked runs, and end \emph{above} the
$1.22$--$1.45$ band that grokked runs occupy early. Independent corroboration from PyTorch runs
at the same modulus reads $1.626$ and $1.574$ at accuracies $0.2236$ and $0.2749$. Two
implementations, two data sources, one answer.

The early ramp is therefore not evidence of circuit formation and must not be read as a
progress measure. This extends rather than contradicts the specificity result: that
already recorded that the permutation $p$ does not separate grokked from failed runs and that
only the enrichment \emph{magnitude} does. What is new is that the ramp's \emph{shape} does not
separate either, and monotonicity was the last property that might have made the early signal
diagnostic. Enrichment magnitude does still separate cleanly:
grokked runs reach $3.3$ to $17.4$ where these failures top out at $1.66$. The ablation
evidence for the CRT structure in grokked models (Section~\ref{sec:causal}) is untouched. Wherever the ramp appears in this
paper, the failed-run ramp appears beside it.

\subsection{Additive sparsity, zero divisors, and $\omega(n)$}
\label{sec:crt-omega}

Additive sparsity tracks zero-divisor density across $18$ moduli at $\rho = +0.858$. It also
separates into non-overlapping bands by $\omega(n)$, the number of distinct prime factors:
$\omega = 1$ spans $0.017$--$0.184$, $\omega = 2$ spans $0.394$--$0.481$, and $\omega = 3$
spans $0.578$--$0.589$, against a residue-axis random-orthogonal control flat at
$0.117$--$0.155$ throughout (Figure~\ref{fig:omega}).

This is a confirmatory result, not a discovery, and the distinction is not cosmetic.
The bands were seen before the pre-registration was written; a pre-registration cannot make an
observation on the same data into a prediction. What the confirmatory test does is kill the
confounds, which is worth doing because three claims in this project were lost to exactly one:
a grouping variable coupled to a size, with the statistic reading the size. Here the
grouping variable $\omega$ is collinear with both zero-divisor density and $n$ itself. The
size confound is absent: a blind test reads $\rho(n, \mathrm{Gini}_{\text{add}}) = -0.185$, the
random-orthogonal control has no band structure at all, and a $10{,}000$-shuffle permutation of
$\omega$ puts the minimum adjacent gap at $p = 1.0 \times 10^{-4}$. At matched zero-divisor
density, $10$ discordant pairs favour $\omega$ $10$ out of $10$ (sign test
$p = 1.95 \times 10^{-3}$; the pairs share moduli and are not independent, so the fraction is
the statistic and the $p$ is descriptive).

$\omega$ does not displace zero-divisor density, and we do not crown it. For
predictors $X, Z$ and response $Y$ the partial Spearman correlation
\begin{equation}
  \label{eq:partial}
  \rho(X, Y \mid Z) \;=\;
  \frac{\rho_{XY} - \rho_{XZ}\,\rho_{ZY}}{\sqrt{(1 - \rho_{XZ}^{2})(1 - \rho_{ZY}^{2})}}
\end{equation}
asks what $X$ explains that $Z$ does not. Both survive it:
$\rho(\omega, G \mid \mathrm{zdd}) = +0.804$ and
$\rho(\mathrm{zdd}, G \mid \omega) = +0.578$, so neither is mediated by the other. We decline to
read an ordering off that gap. Ranks here must be \emph{midranks}. $\omega$ takes three distinct
values across the $23$ moduli, and under the ordinal ranking we used until 2026-09-23 the two
partials read $+0.627$ and $+0.735$. Those are the same two numbers in the opposite order,
because within an $\omega$ band an ordinal rank is the ordering of $n$. A quantity whose sign of
difference depends on a tie convention is not a ranking. The band structure was also seen before
the pre-registration, the matched pairs share moduli and are not independent, and the designed
falsifier at $n = 128$ was ambiguous. The defensible statement is that $\omega(n)$ and
zero-divisor density are two partially independent predictors of additive sparsity and this
data cannot rank them.

The designed falsifier fired and neither predictor was accurate. $n = 128 = 2^{7}$ is
a prime power ($\omega = 1$) with half its residues zero divisors, the one place the two
predictors disagree, and the two were pre-registered as disjoint intervals: $\omega$ predicted
$0.017$--$0.184$ and zero-divisor density predicted $0.434$--$0.589$. It read $\mathbf{0.261}$,
between both. We report this as ambiguous. What it settles is smaller than what it was meant
to settle. $128$ lands on $\omega$'s side and extends the $\omega = 1$ band to
$0.015$--$0.261$, while the $\omega = 1$ and $\omega = 2$ bands still do not overlap
($0.261$ against $0.394$). It also favours $\omega$ in all three of the matched pairs it
appears in.

% STE descriptive
\begin{figure}[t]
  \centering
  \includegraphics[width=\textwidth]{../figures/omega_bands.pdf}
  \caption{\textbf{Additive sparsity by $\omega(n)$, with its confound and its falsifier.}
    The panels give the bands, the zero-divisor-density confound, the flat random-orthogonal
    size control, and the $n = 2^{7}$ prediction drawn \emph{beside} its outcome. That
    outcome landed between the two predicted intervals, so neither predictor was accurate.
    The figure regenerates from \texttt{src/viz/plots.py::omega\_bands}.}
  \label{fig:omega}
\end{figure}
% STE end

$\omega$ is the mechanistically expected variable rather than a fished one. The CRT-dual
set of Section~\ref{sec:crt-law} is built as one family per maximal prime power, so at
$\omega = 1$ that set is every frequency, and there is nothing to concentrate on. Each further prime adds
a family and hence more concentration. It is also why a prime power's clock must live in the
multiplicative basis: additively, there is nothing there.
